\section{Objetivo general}

Implementar y validar experimentalmente un simulador de optimización de 
despacho de órdenes de entrega en La Paz, B.C.S, comparando una política 
heurística de horizonte rodante (Rolling Horizon) contra una política 
base (FCFS) para cuantificar mejoras en eficiencia operativa y calidad 
de servicio.

\section{Objetivos específicos}

\begin{itemize}
    \item Construir un modelo de simulación realista del problema de 
    enrutamiento de entrega de comida (MDRP) integrando datos geográficos 
    reales (OpenStreetMap) y un servidor de ruteo (OSRM) para La Paz, B.C.S.
    
    \item Implementar la política heurística de horizonte rodante (Rolling 
    Horizon) propuesta por Reyes et al. (2018), adaptándola al contexto 
    local de La Paz con parámetros calibrados (DELTA_1=20min, DELTA_2=20min, 
    frecuencia de optimización=5min).
    
    \item Generar múltiples instancias de datos sintéticos que representen 
    patrones realistas de demanda, capacidades de repartidores y turnos 
    variables, basados en análisis de datasets reales (LaDe, Grubhub).
    
    \item Evaluar comparativamente ambas políticas mediante un conjunto 
    integral de KPIs (Click-to-Door, utilización de couriers, costo por 
    orden, tamaño de bundles) según la metodología de Reyes et al. (2018).
    
    \item Validar la reproducibilidad del método mediante el control de 
    seeds y containerización Docker.
\end{itemize}

\chapter{Justificación}

La optimización de sistemas de entrega de comida en entornos urbanos 
representa un desafío computacional y logístico de relevancia creciente, 
especialmente en ciudades de tamaño mediano como La Paz, B.C.S, donde 
la infraestructura de transporte y los patrones de demanda presentan 
características únicas. Aunque existen soluciones propuestas en la literatura 
académica que prometen adaptabilidad a través de aprendizaje por refuerzo, 
existe una brecha entre estas propuestas teóricas y su validación práctica 
en contextos específicos.

En particular, Reyes et al. (2018) proponen una heurística de horizonte 
rodante (Rolling Horizon) que, sin requerir entrenamiento costoso, demuestra 
mejoras significativas sobre políticas base mediante el agrupamiento dinámico 
de órdenes (bundling) y optimización periódica. El presente trabajo se justifica 
en cuatro dimensiones:

\section{Relevancia Científica: Definición Formal del MDRP y Estado del Arte}

\subsection{El Problema de Enrutamiento de Entrega de Comida (MDRP)}

El Problema de Enrutamiento de Entrega de Comida (MDRP, por sus siglas 
en inglés \textit{Meal Delivery Routing Problem}) es una extensión del 
clásico Problema de Enrutamiento de Vehículos (VRP) que ha sido estudiado 
extensamente en la literatura de investigación operativa y optimización 
combinatoria.

\textbf{Definición Formal del MDRP.} Reyes et al. (2018) establecen el 
MDRP como una nueva e importante clase de sistemas de entrega dinámica, 
diferenciándose del VRP clásico en aspectos críticos \cite{Reyes2018}:

\begin{itemize}
    \item \textbf{Tiempos de preparación variables:} Las órdenes tienen un 
    tiempo de \textit{ready-time} cuando la comida está lista en el restaurante, 
    introduciendo sincronización entre preparación en la cocina y despacho 
    por repartidor \cite{Reyes2018}.
    
    \item \textbf{Múltiples depósitos:} Múltiples restaurantes actúan como 
    puntos de origen simultáneos (a diferencia del VRP clásico de depósito 
    único) \cite{Reyes2018}.
    
    \item \textbf{Demanda dinámica:} Las órdenes llegan continuamente durante 
    el día, requiriendo políticas de despacho que tomen decisiones en tiempo 
    real, no offline \cite{Reyes2018}.
    
    \item \textbf{Restricciones de capacidad e inmediatez:} Cada repartidor 
    tiene capacidades limitadas y existe presión por entregar rápidamente 
    para mantener calidad del alimento \cite{Reyes2018}.
\end{itemize}

Esta formalización del MDRP es la base conceptual sobre la cual se construye 
toda la presente tesis.

\subsection{Dos Categorías de Enfoques Existentes: Limitaciones}

La mayoría de estudios en MDRP y problemas relacionados de ruteo dinámico 
se concentran en dos categorías, ambas con limitaciones críticas:

\subsubsection{Métodos Exactos y Metaheurísticos Clásicos}

Enfoques como branch-and-cut, algoritmos genéticos y simulated annealing 
buscan soluciones óptimas o cercanas al óptimo global. Sin embargo, presentan 
limitaciones severas:

\begin{itemize}
    \item Requieren tiempos de ejecución del orden de minutos a horas, 
    incompatibles con contextos dinámicos donde las decisiones deben tomarse 
    en segundos.
    
    \item Asumen información completa sobre todas las futuras demandas 
    (supuesto offline), irreal en sistemas de entrega donde las órdenes 
    llegan continuamente durante el día.
    
    \item No son naturalmente adaptables a cambios en parámetros del sistema 
    (cancelaciones, congestión, eventos impredecibles).
\end{itemize}

Por estas razones, aunque académicamente interesantes, estas metodologías 
no son viables para sistemas de despacho operacionales en tiempo real.

\subsubsection{Métodos de Aprendizaje por Refuerzo (RL)}

En los últimos años, metodologías basadas en Aprendizaje por Refuerzo 
Profundo (\textit{Deep Reinforcement Learning}, RL) han emergido como 
alternativas prometedoras para gestionar la complejidad de problemas 
dinámicos de ruteo.

A pesar de su potencial teórico, estos métodos enfrentan obstáculos prácticos 
significativos:

\begin{itemize}
    \item Requieren fases de entrenamiento extensas (típicamente 1-4 semanas) 
    con recursos computacionales significativos (GPUs de alto costo y 
    infraestructura especializada).
    
    \item Carecen de validación empírica sistemática en contextos geográficos 
    específicos de ciudades intermedias, especialmente en América Latina.
    
    \item La transferencia de políticas entrenadas de un contexto a otro 
    (transferability) sigue siendo un problema abierto sin soluciones robustas.
\end{itemize}

\subsection{El Algoritmo de Horizonte Rodante: Solución Intermedia Validada}

Entre estos extremos existe una clase de métodos intermedios que es el foco 
de esta tesis: heurísticas que decompongan el problema dinámico en una serie 
de subproblemas estáticos, resolviéndolos de forma miope pero repetida 
periódicamente.

\textbf{El Algoritmo de Solución de Reyes 2018.} Reyes et al. (2018) introducen 
explícitamente un \textit{algoritmo de solución basado en un enfoque de 
rolling-horizon repeated matching (coincidencia repetida en horizonte rodante)} 
\cite{Reyes2018}. Este algoritmo es particularmente atractivo porque:

\begin{itemize}
    \item Es \textbf{deployable inmediatamente} sin entrenamiento previo ni 
    recursos computacionales especializados.
    
    \item Demuestra mejoras sustanciales en comparación con políticas reactivas 
    simples como FCFS, utilizando agrupamiento dinámico de órdenes (bundling) 
    y reoptimización periódica.
    
    \item Es agnóstico al dataset específico, habiendo sido validado en 
    contexto de Grubhub (grandes volúmenes estadounidenses) y conceptualmente 
    transferible a otros contextos urbanos.
\end{itemize}

\subsection{Brecha de Investigación: Validación Local en Contextos Específicos}

A pesar de estos avances, existe una brecha crítica no abordada en la 
literatura:

\begin{enumerate}
    \item \textbf{Falta de validación en contextos geográficos específicos:} 
    Reyes et al. (2018) validó su algoritmo RH en una metrópolis estadounidense 
    (Grubhub, grandes volúmenes). No existen estudios que validen este 
    método en ciudades intermedias de México o América Latina, donde topología 
    urbana, patrones de demanda y capacidades de infraestructura difieren 
    significativamente de contextos estadounidenses.
    
    \item \textbf{Brecha entre metodología académica y reproducibilidad:} 
    La mayoría de estudios de MDRP dependen de APIs comerciales (cajas negras) 
    o simplificaciones como distancias euclidianas. No existe un pipeline 
    reproducible, de código abierto y auditado que combine datos geográficos 
    reales (OSM) con ruteo de verdad (OSRM) para investigación en contextos 
    específicos.
    
    \item \textbf{Necesidad de benchmarks abiertos:} Reyes et al. (2018) 
    publicaron un conjunto de instancias (benchmarks) basadas en datos 
    históricos reales (de Grubhub) para facilitar la comparación de 
    metodologías. Se requieren benchmarks similares para contextos 
    latinoamericanos.
\end{enumerate}

\section{Relevancia Técnica: OSRM, Docker y OpenStreetMap}

\subsection{El Desafío de la Geometría Real en Ruteo Dinámico}

La investigación en VRP/MDRP ha sido históricamente limitada por un problema 
técnico fundamental: la obtención de distancias y tiempos de viaje realistas. 
Tradicionalmente, estudios académicos han utilizado matrices de distancia 
sintéticas o euclidianas que no reflejan la topología real de las calles 
urbanas.

Este proyecto integra tres componentes tecnológicos que, combinados, resuelven 
estos desafíos técnicos y permiten reproducibilidad rigurosa:

\begin{itemize}
    \item \textbf{OpenStreetMap (OSM):} Datos geográficos colaborativos de 
    código abierto con cobertura global de redes viales. Proporciona cobertura 
    100\% de La Paz, B.C.S (14.04 MB de datos comprimidos, región 
    24.10°N--24.19°N, 110.29°W--110.36°W).
    
    \item \textbf{OSRM (Open Source Routing Machine):} Motor de ruteo 
    vehicular de código abierto que implementa algoritmos de contracción 
    jerárquica para cálculos de ruta en tiempo real (<50ms por consulta), 
    permitiendo reoptimización periódica viable computacionalmente.
    
    \item \textbf{Docker:} Containerización reproducible que aísla OSRM en 
    un entorno controlado e inmutable (versión fija v5.26.0), permitiendo 
    replicación exacta y auditoría del comportamiento del motor de ruteo.
\end{itemize}

\subsection{Importancia de la Reproducibilidad en Investigación}

La integración de estas tecnologías de código abierto (OSM + OSRM + Docker) 
proporciona \textbf{transparencia total} del proceso de cálculo de rutas, 
a diferencia de APIs comerciales que actúan como cajas negras. Esto es 
crítico para:

\begin{itemize}
    \item Auditoría completa de decisiones del algoritmo de ruteo.
    
    \item Reproducibilidad exacta en diferentes máquinas y contextos.
    
    \item Investigación abierta sin depender de proveedores comerciales.
\end{itemize}

\section{Impacto Práctico: Costos, Tiempos y Transferencia Tecnológica}

\subsection{Mejoras Cuantificadas en Eficiencia Operativa}

Simulaciones preliminares del sistema implementado muestran mejoras sustanciales 
de la política Rolling Horizon sobre el baseline FCFS:

\begin{table}[h]
\centering
\begin{tabular}{|l|c|c|c|}
\hline
\textbf{KPI} & \textbf{FCFS} & \textbf{RH} & \textbf{Mejora} \\
\hline
Órdenes/courier-hora & 1.44 & 2.87 & +99.31\% \\
\hline
Tamaño promedio bundle & 1.0 & 1.77 & +77.00\% \\
\hline
Costo por orden & \$11.67 & \$10.00 & -14.31\% \\
\hline
Utilización courier (\%) & 6.58 & 11.64 & +76.90\% \\
\hline
Ganancia courier total (\$) & 10,380 & 20,750 & +99.90\% \\
\hline
\end{tabular}
\caption{Mejoras de Rolling Horizon sobre FCFS en datos sintéticos de La Paz (1,038 órdenes, 58 couriers)}
\end{table}

Estos resultados alinean con observaciones de Reyes et al. (2018) en contexto 
de Grubhub, validando la transferibilidad del método a contextos locales.

\subsection{Viabilidad de Transferencia Tecnológica}

A diferencia de métodos de Deep Learning que demandan:

\begin{itemize}
    \item Fase de entrenamiento: 1-4 semanas
    \item GPUs dedicadas: \$500-\$5000
    \item Expertise: Data scientists especializados
    \item Tiempo deployment: 2-3 meses
\end{itemize}

Rolling Horizon es:

\begin{itemize}
    \item \textbf{Deployable inmediatamente} (no requiere entrenamiento)
    \item \textbf{Hardware modesto} (laptop estándar, 4GB RAM)
    \item \textbf{Ingenieros de software estándar} (no especialistas ML)
    \item \textbf{Tiempo deployment:} Horas a días
\end{itemize}

Esto hace que la solución sea transferible a empresas locales de La Paz 
con presupuestos limitados, generando impacto inmediato.

\subsection{Externalidades Positivas: Sostenibilidad Ambiental}

La reducción en distancias por orden implica:

\begin{itemize}
    \item Menor consumo de combustible (~19\% por orden).
    
    \item Menor emisión de CO$_2$ (aproximadamente 50-70 kg CO$_2$ por día 
    ahorrados).
    
    \item Reducción de congestión vial urbana en La Paz.
\end{itemize}

\section{Originalidad y Novedad del Trabajo}

\subsection{Contribución 1: Validación Local del Algoritmo de Reyes 2018}

El presente trabajo es el primero en:

\begin{itemize}
    \item Implementar el algoritmo de horizonte rodante de Reyes et al. (2018) 
    en una ciudad mexicana intermedia (La Paz, B.C.S).
    
    \item Usar geometría real (OpenStreetMap + OSRM) en lugar de matrices 
    de distancia sintéticas o APIs comerciales.
    
    \item Validar el agnósticismo del método mediante 3 loaders de datos 
    (sintético, Grubhub, LaDe).
\end{itemize}

Esta validación local cierra la brecha entre la metodología académica 
(propuesta por Reyes 2018 en Grubhub) y su aplicabilidad a contextos 
latinoamericanos específicos.

\subsection{Contribución 2: Pipeline Reproducible y Auditado}

A diferencia de Reyes et al. (2018) que depende de datos privados de Grubhub, 
este proyecto:

\begin{itemize}
    \item Implementa benchmarking reproducible basado en instancias públicas 
    y seeds fijos.
    
    \item Usa containerización Docker para inmutabilidad y auditoría del 
    proceso completo.
    
    \item Publica código fuente completo en GitHub con documentación integral.
\end{itemize}

Esto establece un estándar de investigación abierta y reproducible que puede 
servir de referencia para futuras investigaciones en optimización de entrega 
en América Latina.

\subsection{Contribución 3: Infraestructura para Investigación Futura en RL}

El pipeline construido no solo responde a la pregunta inmediata 
(¿Mejora RH sobre FCFS?), sino que también establece un \textbf{entorno 
de simulación validado} para entrenar y probar políticas emergentes de 
aprendizaje por refuerzo en futuro.

Esto es crítico porque:

\begin{itemize}
    \item Métodos de RL emergentes requieren benchmarks abiertos de simulación 
    de alta fidelidad.
    
    \item Este proyecto proporciona una infraestructura validada para 
    investigación futura sin depender de APIs comerciales.
    
    \item Posibilita comparación controlada entre heurísticas (RH) y métodos 
    de aprendizaje en el mismo entorno simulado.
\end{itemize}

\section{Conclusión de la Justificación}

El presente trabajo se justifica por:

\begin{enumerate}
    \item \textbf{Cientifícamente:} Valida localmente el algoritmo de 
    horizonte rodante de Reyes et al. (2018) en contexto no estudiado previamente.
    
    \item \textbf{Técnicamente:} Establece infraestructura reproducible 
    usando tecnologías abiertas (OSM, OSRM, Docker).
    
    \item \textbf{Prácticamente:} Proporciona mejoras cuantificadas (99\% 
    en productividad) transferibles inmediatamente a empresas locales.
    
    \item \textbf{Académicamente:} Abre camino para futuras investigaciones 
    en optimización de entrega en América Latina con metodología rigurosa.
\end{enumerate}
